\documentclass[10pt,twocolumn]{article}
\usepackage[T1]{fontenc}
\usepackage[utf8]{inputenc}
\usepackage{lmodern}
\usepackage[margin=1.8cm,top=2.4cm,bottom=2cm,columnsep=0.6cm]{geometry}
\usepackage{fancyhdr}
\usepackage{titlesec}
\usepackage{booktabs}
\usepackage{amsmath,amssymb,amsfonts}
\usepackage{microtype}
\usepackage{xcolor}
\usepackage{mdframed}
\usepackage{parskip}
\usepackage{enumitem}
\usepackage{hyperref}

\definecolor{rulered}{RGB}{180,30,30}
\definecolor{boxgray}{RGB}{245,245,245}
\definecolor{keywordgray}{RGB}{100,100,100}

\pagestyle{fancy}
\fancyhf{}
\fancyhead[L]{\textsc{\small Claude Mag}}
\fancyhead[C]{\small\textit{Machine Learning Research Digest}}
\fancyhead[R]{\small 2026-08-12}
\fancyfoot[C]{\small\thepage}
\renewcommand{\headrulewidth}{0.8pt}

\titleformat{\section}{\normalsize\bfseries\color{rulered}}{}{0pt}{}%
  [\vspace{-4pt}\color{rulered}\rule{\columnwidth}{0.4pt}]
\titlespacing{\section}{0pt}{8pt}{4pt}

\hypersetup{colorlinks=true,urlcolor=rulered,linkcolor=black}

\begin{document}
\setlength{\parindent}{0pt}
\setlength{\parskip}{4pt}

{\small\color{keywordgray}\textbf{Keywords:}
  context reliability \textbullet{}
  sycophancy \textbullet{}
  direct preference optimization \textbullet{}
  LLM alignment}\par\vspace{4pt}

{\LARGE\bfseries\raggedright
  Learning When to Trust via Selective Context
  Preference Optimization}\par\vspace{4pt}

{\large\itshape\raggedright
  MIST Benchmark and SCOPE Training Teach Language Models
  to Resist Misleading Context Without Ignoring Helpful Signals}\par\vspace{6pt}

{\small\textbf{By} Xian Sun, Wei Chow, Yingshuo Wang, Junhao Liu,
  Wei Gao, Qing Wu, Lingdong Kong}\par\vspace{2pt}

{\small\color{keywordgray}arXiv:2608.06377 \textbullet{} Preprint, August 2026}
\par\vspace{2pt}\noindent{\color{rulered}\rule{\columnwidth}{0.6pt}}\vspace{6pt}

A new paper introduces MIST, a benchmark that renders each reasoning item under
four matched context conditions, and SCOPE, a DPO-based training method that teaches
models to distinguish misleading from helpful context. SCOPE significantly reduces
the fraction of correct answers flipped wrong by misleading signals while preserving
the model's benefit from genuinely useful context.

\begin{mdframed}[backgroundcolor=boxgray,linecolor=rulered,linewidth=1pt,
    innerleftmargin=6pt,innerrightmargin=6pt,
    innertopmargin=6pt,innerbottommargin=6pt]
\textbf{\small AT A GLANCE}\par\vspace{4pt}
\begin{description}[leftmargin=0pt,labelsep=4pt,itemsep=2pt,parsep=0pt,
    font=\small\bfseries]
  \item[Problem:] LLMs are sycophantic toward misleading context yet training to
    resist context can make them ignore helpful signals.
  \item[Why it matters:] RAG, web search, and chat history all deliver context
    that may be wrong; models that cannot discriminate are unsafe to deploy.
  \item[Method:] MIST benchmark (4 matched conditions per item) + SC2W metric +
    SCOPE (DPO over balanced preference pairs across all four conditions).
  \item[Result:] SCOPE significantly reduces SC2W while preserving correct-context
    utilization; balanced training is crucial.
  \item[Evidence:] Ablation confirms that optimizing only on misleading pairs
    degrades helpful-context performance.
\end{description}
\end{mdframed}

\section{The Big Picture}

Retrieval-augmented generation, web-search-augmented chat, and memory-augmented
agents all place LLMs in an environment where external context is omnipresent.
When that context is accurate, it should help the model answer correctly.
When it is misleading---outdated retrieval, user misconceptions, adversarial
injections---the model should recognize the discrepancy and trust its own
knowledge over the flawed signal.

Current LLMs fail in both directions: they are sycophantic toward user-provided
context (a misleading signal flips a correct answer wrong), yet attempts to harden
models against misleading context often produce models that ignore helpful context
universally. This paper frames the problem as a \textbf{discrimination task}: the
model must learn to distinguish the quality of context, not merely resist all of it.

\section{Why Existing Approaches Fall Short}

Prior robustness work has largely attacked misleading context as a single phenomenon,
training on adversarial examples where context is always wrong. This produces models
with one failure mode traded for another: high resistance to misleading context, but
equally high resistance to correct context, because the model has learned
``context $\to$ ignore,'' not ``misleading context $\to$ ignore.''

Evaluating this nuance requires a benchmark that explicitly presents both misleading
and correct context for the same item. Without matched conditions, it is impossible
to separate context-resistance from context-utilization, and prior evaluations have
conflated the two.

\section{Core Mechanism}

\textbf{MIST Benchmark.}
Each item in MIST appears under four matched conditions:
\begin{enumerate}[itemsep=2pt,parsep=0pt]
  \item \textbf{Clean:} question alone, no external context
  \item \textbf{Misleading:} context that nudges toward the wrong answer
  \item \textbf{Correct context:} context that supports the right answer
  \item \textbf{Irrelevant:} context unrelated to the question
\end{enumerate}

The \textbf{SC2W metric} (Susceptibility to Context-induced Wrong answers) counts
how often condition~2 flips a clean-correct answer to wrong:
\[
\mathrm{SC2W} = \frac{\#\{\text{clean-correct} \wedge \text{misleading-wrong}\}}{
  \#\{\text{clean-correct}\}}
\]

\textbf{SCOPE Training.}
SCOPE (Selective Context Preference Optimization) applies DPO over preference pairs
mined from MIST failures. The preferred completion is the model's clean-correct
answer; the dispreferred completion is the misleading-flipped wrong answer.
Crucially, pairs are sampled \textbf{equally across all four conditions},
not only misleading ones.

\section{Technical Formulation}

Let $x$ be a question and $c$ a context drawn from condition set
$\mathcal{C} = \{\emptyset, \text{mis}, \text{cor}, \text{irr}\}$.
The DPO objective over preference pairs $(y_w, y_l)$ is:

\[
\mathcal{L}_{\mathrm{DPO}} = -\mathbb{E}_{(x,c,y_w,y_l) \sim \mathcal{D}}\!\left[
  \log\sigma\!\left(\beta\log\frac{\pi_\theta(y_w \mid x,c)}{\pi_{\mathrm{ref}}(y_w \mid x,c)}
  - \beta\log\frac{\pi_\theta(y_l \mid x,c)}{\pi_{\mathrm{ref}}(y_l \mid x,c)}\right)\right]
\]

where $\mathcal{D}$ is the balanced dataset spanning all four conditions with equal weight.
The SC2W constraint requires that:
\[
\mathbb{P}\bigl[y^* \mid x, c_{\mathrm{mis}}\bigr] >
\mathbb{P}\bigl[y^* \mid x, c_{\mathrm{mis}}\bigr]\big|_{\pi_{\mathrm{ref}}}
\]
while preserving:
\[
\mathbb{P}\bigl[y^* \mid x, c_{\mathrm{cor}}\bigr] \approx
\mathbb{P}\bigl[y^* \mid x, c_{\mathrm{cor}}\bigr]\big|_{\pi_{\mathrm{ref}}}
\]

\section{Learning or Inference Procedure}

\textbf{Pair mining:} Identify items where the base model is correct in the clean
condition but incorrect in the misleading condition. Each such item yields one DPO
preference pair.

\textbf{Balanced sampling:} Sample equal numbers of preference pairs from each of
the four conditions. This ensures the model cannot learn a context-agnostic
``ignore all signals'' policy.

\textbf{DPO fine-tuning:} Standard DPO applied to the preference dataset; no
modifications to the DPO algorithm are required.

\textbf{Inference:} No change. The trained model processes queries with external
context in the standard RAG pipeline.

\section{What the Guarantee Says}

No formal guarantee. The balanced training design provides an intuitive argument:
because pairs from the correct-context condition reward trusting helpful context
while pairs from the misleading condition reward resisting bad context, the model
must learn the discrimination rather than a blanket policy. The empirical ablation
confirms this: training on misleading-only pairs degrades correct-context
utilization, while balanced training does not.

\section{Experimental Findings}

\begin{itemize}[itemsep=2pt,parsep=0pt]
  \item \textbf{SC2W reduction:} SCOPE significantly reduces the fraction of
    correct answers flipped wrong by misleading context.
  \item \textbf{Correct-context utilization:} Preserved after SCOPE training;
    the model continues to benefit from genuine context.
  \item \textbf{Irrelevant-context:} Model correctly learns to ignore
    irrelevant signals.
  \item \textbf{Clean accuracy:} Maintained; SCOPE does not degrade the model's
    baseline reasoning.
\end{itemize}

\section{Ablations and Interpretation}

\textbf{Misleading-only DPO vs.\ SCOPE.}
Training DPO only on misleading-context pairs achieves strong SC2W reduction
but degrades correct-context utilization by a significant margin, confirming
the ``context-ignore'' failure mode. SCOPE's balanced sampling resolves this.

\textbf{Effect of condition balance.}
Increasing the proportion of misleading pairs toward 100\% monotonically
trades SC2W improvement for correct-context degradation. The balanced
50/50/50/50 split across conditions achieves the best combined performance.

\textbf{MIST item difficulty.}
High-SC2W items tend to involve superficially plausible misleading contexts
(wrong dates, plausible-but-incorrect numbers). These are the cases where
the model most needs the discrimination capability SCOPE provides.

\section*{Reference}

Xian Sun, Wei Chow, Yingshuo Wang, Junhao Liu, Wei Gao, Qing Wu, Lingdong Kong.
\textbf{Learning When to Trust via Selective Context Preference Optimization}.
arXiv:2608.06377, August 2026.
\url{https://arxiv.org/abs/2608.06377}

\end{document}
